% --- Archon physics formalization source begin ---
% archon:physics
% archon:covers PhyXMiniProblems/problem_phyx_mini_0444.lean
% archon:source-report reports/phyx_mini/problem_phyx_mini_0444.source.json

\chapter{Physics problem phyx\_mini\_0444}
\label{ch:PhyXMiniProblems_problem_phyx_mini_0444}

\paragraph{Problem source.}
\#\# Physical scenario

Consider two tanks, \$A\$ and \$B\$, connected by a valve, as shown in figure. Each has a volume of \$200 \textbackslash{}, \textbackslash{}text\{L\}\$, and tank \$A\$ has R-410a at \$25 \textbackslash{}, \textasciicircum{}\{\textbackslash{}circ\}\textbackslash{}text\{C\}\$, \$10\textbackslash{}\%\$ liquid and \$90\textbackslash{}\%\$ vapor by volume, while tank \$B\$ is evacuated. The valve is now opened, and saturated vapor flows from \$A\$ to \$B\$ until the pressure in \$B\$ has reached that in \$A\$, at which point the valve is closed. This process occurs slowly such that all temperatures stay at \$25 \textbackslash{}, \textasciicircum{}\{\textbackslash{}circ\}	ext\{C\}\$ throughout the process.

\#\# Question

How much has the quality changed in tank \$A\$ during the process?

\#\# Answer choices

- A: A: 0.016

- B: B: 0.500

- C: C: 0.159

- D: D: 0.283

\#\# Figure caption (auxiliary; use the image as primary evidence)

The image shows two containers labeled "A" and "B." Container A has some liquid at the bottom, represented with a blue fill and horizontal lines indicating water. Container B is empty. Both containers are connected to a light bulb represented in the top center by two lines leading to it from the containers. The light bulb is depicted by a circle with an "X" inside it. This setup suggests a simple circuit or illustration of electrical or fluid dynamics concepts.

\#\# Dataset metadata

Laws of Thermodynamics; Physical Model Grounding Reasoning, Multi-Formula Reasoning

\paragraph{Recorded answer/context.}
D: D: 0.283

\paragraph{Figure/image path.}
/root/proposal\_for\_physic/science-mango/phyx\_mini\_run/phyx\_data/test\_image/444.png

\paragraph{Formalization target.}
create a compiling Lean file with sorry bodies at `PhyXMiniProblems/problem\_phyx\_mini\_0444.lean`. The Lean declarations must preserve the physical quantities, dimensions or dimensional roles, figure labels, governing-law hypotheses, and final relation expressed by this problem.
Use Mathlib/Physlib names found through LeanExplore where available. If a physics API is missing, introduce faithful local abstractions rather than scalar placeholder aliases.

\begin{theorem}[Physics formalization target]
\label{thm:physics:phyx_mini_0444:target}
\lean{PhyXMiniProblems.ProblemPhyXMini0444.quality_change_in_tank_A_is_choice_D}
\uses{def:physics:phyx-mini-0444:phyxminiproblems-problemphyxmini0444-twotankr410aexperiment, def:physics:phyx-mini-0444:phyxminiproblems-problemphyxmini0444-qualitychangeintanka, def:physics:phyx-mini-0444:phyxminiproblems-problemphyxmini0444-matchessuppliedtwotankfigure, def:physics:phyx-mini-0444:phyxminiproblems-problemphyxmini0444-matchesr410atankscenario, def:physics:phyx-mini-0444:phyxminiproblems-problemphyxmini0444-matchesreferencer410apropertydata, def:physics:phyx-mini-0444:phyxminiproblems-problemphyxmini0444-hasphysicalr410aparameters, def:physics:phyx-mini-0444:phyxminiproblems-problemphyxmini0444-satisfiesisothermalsaturatedtransferlaws, def:physics:phyx-mini-0444:phyxminiproblems-problemphyxmini0444-answerchoice, def:physics:phyx-mini-0444:phyxminiproblems-problemphyxmini0444-isreportedqualitychange}
The assigned autoformalize agent should translate this physics problem into Lean declarations in the covered file, with theorem and lemma proof bodies written as `by sorry`.
\end{theorem}
\begin{proof}
This is an autoformalization task, not a proof task. Produce statements that can later be proved by the physics prover without weakening the physical model.
\end{proof}

% --- Archon declaration topology begin ---
\section{Declaration topology}
\begin{definition}[volume Dimension]
  \label{def:physics:phyx-mini-0444:phyxminiproblems-problemphyxmini0444-volumedimension}
  \lean{PhyXMiniProblems.ProblemPhyXMini0444.volumeDimension}
  The physical dimension of volume, L³
\end{definition}
\begin{proof}
  This is the declared model definition of volume Dimension.
\end{proof}
\begin{definition}[specific Volume Dimension]
  \label{def:physics:phyx-mini-0444:phyxminiproblems-problemphyxmini0444-specificvolumedimension}
  \lean{PhyXMiniProblems.ProblemPhyXMini0444.specificVolumeDimension}
  \uses{def:physics:phyx-mini-0444:phyxminiproblems-problemphyxmini0444-volumedimension}
  The physical dimension of specific volume, L³ M⁻¹
\end{definition}
\begin{proof}
  This is the declared model definition of specific Volume Dimension.
\end{proof}
\begin{definition}[Volume Quantity]
  \label{def:physics:phyx-mini-0444:phyxminiproblems-problemphyxmini0444-volumequantity}
  \lean{PhyXMiniProblems.ProblemPhyXMini0444.VolumeQuantity}
  \uses{def:physics:phyx-mini-0444:phyxminiproblems-problemphyxmini0444-volumedimension}
  A nonnegative physical volume
\end{definition}
\begin{proof}
  This is the declared model definition of Volume Quantity.
\end{proof}
\begin{definition}[Mass Quantity]
  \label{def:physics:phyx-mini-0444:phyxminiproblems-problemphyxmini0444-massquantity}
  \lean{PhyXMiniProblems.ProblemPhyXMini0444.MassQuantity}
  A nonnegative physical mass
\end{definition}
\begin{proof}
  This is the declared model definition of Mass Quantity.
\end{proof}
\begin{definition}[Specific Volume Quantity]
  \label{def:physics:phyx-mini-0444:phyxminiproblems-problemphyxmini0444-specificvolumequantity}
  \lean{PhyXMiniProblems.ProblemPhyXMini0444.SpecificVolumeQuantity}
  \uses{def:physics:phyx-mini-0444:phyxminiproblems-problemphyxmini0444-specificvolumedimension}
  A nonnegative thermodynamic specific volume
\end{definition}
\begin{proof}
  This is the declared model definition of Specific Volume Quantity.
\end{proof}
\begin{definition}[volume In Cubic Metres]
  \label{def:physics:phyx-mini-0444:phyxminiproblems-problemphyxmini0444-volumeincubicmetres}
  \lean{PhyXMiniProblems.ProblemPhyXMini0444.volumeInCubicMetres}
  \uses{def:physics:phyx-mini-0444:phyxminiproblems-problemphyxmini0444-volumequantity}
  Read a physical volume in coherent SI cubic metres
\end{definition}
\begin{proof}
  This is the declared model definition of volume In Cubic Metres.
\end{proof}
\begin{definition}[volume In Litres]
  \label{def:physics:phyx-mini-0444:phyxminiproblems-problemphyxmini0444-volumeinlitres}
  \lean{PhyXMiniProblems.ProblemPhyXMini0444.volumeInLitres}
  \uses{def:physics:phyx-mini-0444:phyxminiproblems-problemphyxmini0444-volumequantity, def:physics:phyx-mini-0444:phyxminiproblems-problemphyxmini0444-volumeincubicmetres}
  Read a physical volume in litres
\end{definition}
\begin{proof}
  This is the declared model definition of volume In Litres.
\end{proof}
\begin{definition}[mass In Kilograms]
  \label{def:physics:phyx-mini-0444:phyxminiproblems-problemphyxmini0444-massinkilograms}
  \lean{PhyXMiniProblems.ProblemPhyXMini0444.massInKilograms}
  \uses{def:physics:phyx-mini-0444:phyxminiproblems-problemphyxmini0444-massquantity}
  Read a physical mass in coherent SI kilograms
\end{definition}
\begin{proof}
  This is the declared model definition of mass In Kilograms.
\end{proof}
\begin{definition}[specific Volume In Cubic Metres Per Kilogram]
  \label{def:physics:phyx-mini-0444:phyxminiproblems-problemphyxmini0444-specificvolumeincubicmetresperkilogram}
  \lean{PhyXMiniProblems.ProblemPhyXMini0444.specificVolumeInCubicMetresPerKilogram}
  \uses{def:physics:phyx-mini-0444:phyxminiproblems-problemphyxmini0444-specificvolumequantity}
  Read a specific volume in cubic metres per kilogram
\end{definition}
\begin{proof}
  This is the declared model definition of specific Volume In Cubic Metres Per Kilogram.
\end{proof}
\begin{definition}[pressure In Pascals]
  \label{def:physics:phyx-mini-0444:phyxminiproblems-problemphyxmini0444-pressureinpascals}
  \lean{PhyXMiniProblems.ProblemPhyXMini0444.pressureInPascals}
  Read a physical pressure in coherent SI pascals
\end{definition}
\begin{proof}
  This is the declared model definition of pressure In Pascals.
\end{proof}
\begin{definition}[Measured Temperature]
  \label{def:physics:phyx-mini-0444:phyxminiproblems-problemphyxmini0444-measuredtemperature}
  \lean{PhyXMiniProblems.ProblemPhyXMini0444.MeasuredTemperature}
  Physlib represents absolute temperature in a zero-preserving unit. A Celsius value is therefore retained as an instrument readout attached to the same absolute temperature, with its affine calibration imposed below
\end{definition}
\begin{proof}
  This is the declared model definition of Measured Temperature.
\end{proof}
\begin{definition}[temperature In Kelvin]
  \label{def:physics:phyx-mini-0444:phyxminiproblems-problemphyxmini0444-temperatureinkelvin}
  \lean{PhyXMiniProblems.ProblemPhyXMini0444.temperatureInKelvin}
  Read an absolute temperature in kelvins
\end{definition}
\begin{proof}
  This is the declared model definition of temperature In Kelvin.
\end{proof}
\begin{definition}[Tank Label]
  \label{def:physics:phyx-mini-0444:phyxminiproblems-problemphyxmini0444-tanklabel}
  \lean{PhyXMiniProblems.ProblemPhyXMini0444.TankLabel}
  The two tank labels printed in the source image
\end{definition}
\begin{proof}
  This is the declared model definition of Tank Label.
\end{proof}
\begin{definition}[System Stage]
  \label{def:physics:phyx-mini-0444:phyxminiproblems-problemphyxmini0444-systemstage}
  \lean{PhyXMiniProblems.ProblemPhyXMini0444.SystemStage}
  The equilibrium endpoints of the transfer process
\end{definition}
\begin{proof}
  This is the declared model definition of System Stage.
\end{proof}
\begin{definition}[Refrigerant Phase]
  \label{def:physics:phyx-mini-0444:phyxminiproblems-problemphyxmini0444-refrigerantphase}
  \lean{PhyXMiniProblems.ProblemPhyXMini0444.RefrigerantPhase}
  Liquid and vapor contributions to a tank inventory
\end{definition}
\begin{proof}
  This is the declared model definition of Refrigerant Phase.
\end{proof}
\begin{definition}[Refrigerant Kind]
  \label{def:physics:phyx-mini-0444:phyxminiproblems-problemphyxmini0444-refrigerantkind}
  \lean{PhyXMiniProblems.ProblemPhyXMini0444.RefrigerantKind}
  Refrigerant identity, kept distinct from scalar property readouts
\end{definition}
\begin{proof}
  This is the declared model definition of Refrigerant Kind.
\end{proof}
\begin{definition}[Tank Contents Region]
  \label{def:physics:phyx-mini-0444:phyxminiproblems-problemphyxmini0444-tankcontentsregion}
  \lean{PhyXMiniProblems.ProblemPhyXMini0444.TankContentsRegion}
  Thermodynamic region occupied by the contents of one tank
\end{definition}
\begin{proof}
  This is the declared model definition of Tank Contents Region.
\end{proof}
\begin{definition}[Valve Position]
  \label{def:physics:phyx-mini-0444:phyxminiproblems-problemphyxmini0444-valveposition}
  \lean{PhyXMiniProblems.ProblemPhyXMini0444.ValvePosition}
  Operational position of the valve
\end{definition}
\begin{proof}
  This is the declared model definition of Valve Position.
\end{proof}
\begin{definition}[Transfer Direction]
  \label{def:physics:phyx-mini-0444:phyxminiproblems-problemphyxmini0444-transferdirection}
  \lean{PhyXMiniProblems.ProblemPhyXMini0444.TransferDirection}
  Direction of net refrigerant transfer through the connecting line
\end{definition}
\begin{proof}
  This is the declared model definition of Transfer Direction.
\end{proof}
\begin{definition}[Transfer Regime]
  \label{def:physics:phyx-mini-0444:phyxminiproblems-problemphyxmini0444-transferregime}
  \lean{PhyXMiniProblems.ProblemPhyXMini0444.TransferRegime}
  Slow-process idealization stated in the prose
\end{definition}
\begin{proof}
  This is the declared model definition of Transfer Regime.
\end{proof}
\begin{definition}[Tank Shape]
  \label{def:physics:phyx-mini-0444:phyxminiproblems-problemphyxmini0444-tankshape}
  \lean{PhyXMiniProblems.ProblemPhyXMini0444.TankShape}
  Rectangular tank outline visible in the supplied raster
\end{definition}
\begin{proof}
  This is the declared model definition of Tank Shape.
\end{proof}
\begin{definition}[Tank State]
  \label{def:physics:phyx-mini-0444:phyxminiproblems-problemphyxmini0444-tankstate}
  \lean{PhyXMiniProblems.ProblemPhyXMini0444.TankState}
  \uses{def:physics:phyx-mini-0444:phyxminiproblems-problemphyxmini0444-volumequantity, def:physics:phyx-mini-0444:phyxminiproblems-problemphyxmini0444-massquantity, def:physics:phyx-mini-0444:phyxminiproblems-problemphyxmini0444-measuredtemperature, def:physics:phyx-mini-0444:phyxminiproblems-problemphyxmini0444-refrigerantphase, def:physics:phyx-mini-0444:phyxminiproblems-problemphyxmini0444-refrigerantkind, def:physics:phyx-mini-0444:phyxminiproblems-problemphyxmini0444-tankcontentsregion}
  Physical state and phase inventory of one tank at one endpoint
\end{definition}
\begin{proof}
  This is the declared model definition of Tank State.
\end{proof}
\begin{definition}[Refrigerant Transfer]
  \label{def:physics:phyx-mini-0444:phyxminiproblems-problemphyxmini0444-refrigeranttransfer}
  \lean{PhyXMiniProblems.ProblemPhyXMini0444.RefrigerantTransfer}
  \uses{def:physics:phyx-mini-0444:phyxminiproblems-problemphyxmini0444-massquantity, def:physics:phyx-mini-0444:phyxminiproblems-problemphyxmini0444-refrigerantphase, def:physics:phyx-mini-0444:phyxminiproblems-problemphyxmini0444-refrigerantkind, def:physics:phyx-mini-0444:phyxminiproblems-problemphyxmini0444-transferdirection}
  The integrated stream that passes through the valve
\end{definition}
\begin{proof}
  This is the declared model definition of Refrigerant Transfer.
\end{proof}
\begin{definition}[Supplied Two Tank Figure]
  \label{def:physics:phyx-mini-0444:phyxminiproblems-problemphyxmini0444-suppliedtwotankfigure}
  \lean{PhyXMiniProblems.ProblemPhyXMini0444.SuppliedTwoTankFigure}
  \uses{def:physics:phyx-mini-0444:phyxminiproblems-problemphyxmini0444-tanklabel, def:physics:phyx-mini-0444:phyxminiproblems-problemphyxmini0444-tankshape}
  Qualitative content of the primary raster. Horizontal coordinates encode only left-to-right placement in the drawing, not physical lengths
\end{definition}
\begin{proof}
  This is the declared model definition of Supplied Two Tank Figure.
\end{proof}
\begin{definition}[R410 AProperty Table]
  \label{def:physics:phyx-mini-0444:phyxminiproblems-problemphyxmini0444-r410apropertytable}
  \lean{PhyXMiniProblems.ProblemPhyXMini0444.R410APropertyTable}
  \uses{def:physics:phyx-mini-0444:phyxminiproblems-problemphyxmini0444-specificvolumequantity, def:physics:phyx-mini-0444:phyxminiproblems-problemphyxmini0444-refrigerantphase}
  An external equilibrium-property table for saturated R-410A. Its outputs are physical pressure and specific volume rather than untyped scalars
\end{definition}
\begin{proof}
  This is the declared model definition of R410 AProperty Table.
\end{proof}
\begin{definition}[Two Tank R410 AExperiment]
  \label{def:physics:phyx-mini-0444:phyxminiproblems-problemphyxmini0444-twotankr410aexperiment}
  \lean{PhyXMiniProblems.ProblemPhyXMini0444.TwoTankR410AExperiment}
  \uses{def:physics:phyx-mini-0444:phyxminiproblems-problemphyxmini0444-volumequantity, def:physics:phyx-mini-0444:phyxminiproblems-problemphyxmini0444-measuredtemperature, def:physics:phyx-mini-0444:phyxminiproblems-problemphyxmini0444-tanklabel, def:physics:phyx-mini-0444:phyxminiproblems-problemphyxmini0444-systemstage, def:physics:phyx-mini-0444:phyxminiproblems-problemphyxmini0444-valveposition, def:physics:phyx-mini-0444:phyxminiproblems-problemphyxmini0444-transferregime, def:physics:phyx-mini-0444:phyxminiproblems-problemphyxmini0444-tankstate, def:physics:phyx-mini-0444:phyxminiproblems-problemphyxmini0444-refrigeranttransfer, def:physics:phyx-mini-0444:phyxminiproblems-problemphyxmini0444-suppliedtwotankfigure, def:physics:phyx-mini-0444:phyxminiproblems-problemphyxmini0444-r410apropertytable}
  The complete experiment. Tank volume is stage-independent, expressing rigidity. The transfer mass and endpoint states are independent observables until the governing laws below are assumed
\end{definition}
\begin{proof}
  This is the declared model definition of Two Tank R410 AExperiment.
\end{proof}
\begin{definition}[phase Mass In Kilograms]
  \label{def:physics:phyx-mini-0444:phyxminiproblems-problemphyxmini0444-phasemassinkilograms}
  \lean{PhyXMiniProblems.ProblemPhyXMini0444.phaseMassInKilograms}
  \uses{def:physics:phyx-mini-0444:phyxminiproblems-problemphyxmini0444-massinkilograms, def:physics:phyx-mini-0444:phyxminiproblems-problemphyxmini0444-tanklabel, def:physics:phyx-mini-0444:phyxminiproblems-problemphyxmini0444-systemstage, def:physics:phyx-mini-0444:phyxminiproblems-problemphyxmini0444-refrigerantphase, def:physics:phyx-mini-0444:phyxminiproblems-problemphyxmini0444-twotankr410aexperiment}
  Kilograms of one phase in a specified endpoint tank
\end{definition}
\begin{proof}
  This is the declared model definition of phase Mass In Kilograms.
\end{proof}
\begin{definition}[total Mass In Kilograms]
  \label{def:physics:phyx-mini-0444:phyxminiproblems-problemphyxmini0444-totalmassinkilograms}
  \lean{PhyXMiniProblems.ProblemPhyXMini0444.totalMassInKilograms}
  \uses{def:physics:phyx-mini-0444:phyxminiproblems-problemphyxmini0444-tanklabel, def:physics:phyx-mini-0444:phyxminiproblems-problemphyxmini0444-systemstage, def:physics:phyx-mini-0444:phyxminiproblems-problemphyxmini0444-twotankr410aexperiment, def:physics:phyx-mini-0444:phyxminiproblems-problemphyxmini0444-phasemassinkilograms}
  Total refrigerant mass in one endpoint tank
\end{definition}
\begin{proof}
  This is the declared model definition of total Mass In Kilograms.
\end{proof}
\begin{definition}[phase Volume In Cubic Metres]
  \label{def:physics:phyx-mini-0444:phyxminiproblems-problemphyxmini0444-phasevolumeincubicmetres}
  \lean{PhyXMiniProblems.ProblemPhyXMini0444.phaseVolumeInCubicMetres}
  \uses{def:physics:phyx-mini-0444:phyxminiproblems-problemphyxmini0444-volumeincubicmetres, def:physics:phyx-mini-0444:phyxminiproblems-problemphyxmini0444-tanklabel, def:physics:phyx-mini-0444:phyxminiproblems-problemphyxmini0444-systemstage, def:physics:phyx-mini-0444:phyxminiproblems-problemphyxmini0444-refrigerantphase, def:physics:phyx-mini-0444:phyxminiproblems-problemphyxmini0444-twotankr410aexperiment}
  Cubic metres occupied by one phase in a specified endpoint tank
\end{definition}
\begin{proof}
  This is the declared model definition of phase Volume In Cubic Metres.
\end{proof}
\begin{definition}[total Occupied Volume In Cubic Metres]
  \label{def:physics:phyx-mini-0444:phyxminiproblems-problemphyxmini0444-totaloccupiedvolumeincubicmetres}
  \lean{PhyXMiniProblems.ProblemPhyXMini0444.totalOccupiedVolumeInCubicMetres}
  \uses{def:physics:phyx-mini-0444:phyxminiproblems-problemphyxmini0444-tanklabel, def:physics:phyx-mini-0444:phyxminiproblems-problemphyxmini0444-systemstage, def:physics:phyx-mini-0444:phyxminiproblems-problemphyxmini0444-twotankr410aexperiment, def:physics:phyx-mini-0444:phyxminiproblems-problemphyxmini0444-phasevolumeincubicmetres}
  Total volume occupied by liquid and vapor in one endpoint tank
\end{definition}
\begin{proof}
  This is the declared model definition of total Occupied Volume In Cubic Metres.
\end{proof}
\begin{definition}[vapor Quality]
  \label{def:physics:phyx-mini-0444:phyxminiproblems-problemphyxmini0444-vaporquality}
  \lean{PhyXMiniProblems.ProblemPhyXMini0444.vaporQuality}
  \uses{def:physics:phyx-mini-0444:phyxminiproblems-problemphyxmini0444-tanklabel, def:physics:phyx-mini-0444:phyxminiproblems-problemphyxmini0444-systemstage, def:physics:phyx-mini-0444:phyxminiproblems-problemphyxmini0444-twotankr410aexperiment, def:physics:phyx-mini-0444:phyxminiproblems-problemphyxmini0444-phasemassinkilograms, def:physics:phyx-mini-0444:phyxminiproblems-problemphyxmini0444-totalmassinkilograms}
  Thermodynamic quality is vapor mass divided by total refrigerant mass. This definition contains no answer-choice value
\end{definition}
\begin{proof}
  This is the declared model definition of vapor Quality.
\end{proof}
\begin{definition}[quality Change In Tank A]
  \label{def:physics:phyx-mini-0444:phyxminiproblems-problemphyxmini0444-qualitychangeintanka}
  \lean{PhyXMiniProblems.ProblemPhyXMini0444.qualityChangeInTankA}
  \uses{def:physics:phyx-mini-0444:phyxminiproblems-problemphyxmini0444-twotankr410aexperiment, def:physics:phyx-mini-0444:phyxminiproblems-problemphyxmini0444-vaporquality}
  Final minus initial vapor quality in tank A
\end{definition}
\begin{proof}
  This is the declared model definition of quality Change In Tank A.
\end{proof}
\begin{definition}[Matches Supplied Two Tank Figure]
  \label{def:physics:phyx-mini-0444:phyxminiproblems-problemphyxmini0444-matchessuppliedtwotankfigure}
  \lean{PhyXMiniProblems.ProblemPhyXMini0444.MatchesSuppliedTwoTankFigure}
  \uses{def:physics:phyx-mini-0444:phyxminiproblems-problemphyxmini0444-suppliedtwotankfigure}
  Facts visible in the primary two-tank raster
\end{definition}
\begin{proof}
  This is the declared model definition of Matches Supplied Two Tank Figure.
\end{proof}
\begin{definition}[Matches R410 ATank Scenario]
  \label{def:physics:phyx-mini-0444:phyxminiproblems-problemphyxmini0444-matchesr410atankscenario}
  \lean{PhyXMiniProblems.ProblemPhyXMini0444.MatchesR410ATankScenario}
  \uses{def:physics:phyx-mini-0444:phyxminiproblems-problemphyxmini0444-volumeinlitres, def:physics:phyx-mini-0444:phyxminiproblems-problemphyxmini0444-pressureinpascals, def:physics:phyx-mini-0444:phyxminiproblems-problemphyxmini0444-twotankr410aexperiment, def:physics:phyx-mini-0444:phyxminiproblems-problemphyxmini0444-phasemassinkilograms, def:physics:phyx-mini-0444:phyxminiproblems-problemphyxmini0444-phasevolumeincubicmetres}
  Prose data and qualitative process facts. The final pressure equality is the stated stopping condition. No field fixes either endpoint quality or its change
\end{definition}
\begin{proof}
  This is the declared model definition of Matches R410 ATank Scenario.
\end{proof}
\begin{definition}[Matches Reference R410 AProperty Data]
  \label{def:physics:phyx-mini-0444:phyxminiproblems-problemphyxmini0444-matchesreferencer410apropertydata}
  \lean{PhyXMiniProblems.ProblemPhyXMini0444.MatchesReferenceR410APropertyData}
  \uses{def:physics:phyx-mini-0444:phyxminiproblems-problemphyxmini0444-specificvolumeincubicmetresperkilogram, def:physics:phyx-mini-0444:phyxminiproblems-problemphyxmini0444-twotankr410aexperiment}
  Calibrated saturation-table readouts at 25 °C. The rounded values v\_f = 0.0009444 m³/kg and v\_g = 0.015164 m³/kg are constitutive R-410A property data, not the requested quality change
\end{definition}
\begin{proof}
  This is the declared model definition of Matches Reference R410 AProperty Data.
\end{proof}
\begin{definition}[Has Physical R410 AParameters]
  \label{def:physics:phyx-mini-0444:phyxminiproblems-problemphyxmini0444-hasphysicalr410aparameters}
  \lean{PhyXMiniProblems.ProblemPhyXMini0444.HasPhysicalR410AParameters}
  \uses{def:physics:phyx-mini-0444:phyxminiproblems-problemphyxmini0444-volumeincubicmetres, def:physics:phyx-mini-0444:phyxminiproblems-problemphyxmini0444-massinkilograms, def:physics:phyx-mini-0444:phyxminiproblems-problemphyxmini0444-specificvolumeincubicmetresperkilogram, def:physics:phyx-mini-0444:phyxminiproblems-problemphyxmini0444-temperatureinkelvin, def:physics:phyx-mini-0444:phyxminiproblems-problemphyxmini0444-twotankr410aexperiment, def:physics:phyx-mini-0444:phyxminiproblems-problemphyxmini0444-phasemassinkilograms, def:physics:phyx-mini-0444:phyxminiproblems-problemphyxmini0444-totalmassinkilograms, def:physics:phyx-mini-0444:phyxminiproblems-problemphyxmini0444-phasevolumeincubicmetres}
  Positivity and nondegeneracy of the physical state parameters
\end{definition}
\begin{proof}
  This is the declared model definition of Has Physical R410 AParameters.
\end{proof}
\begin{definition}[Satisfies Isothermal Saturated Transfer Laws]
  \label{def:physics:phyx-mini-0444:phyxminiproblems-problemphyxmini0444-satisfiesisothermalsaturatedtransferlaws}
  \lean{PhyXMiniProblems.ProblemPhyXMini0444.SatisfiesIsothermalSaturatedTransferLaws}
  \uses{def:physics:phyx-mini-0444:phyxminiproblems-problemphyxmini0444-volumeincubicmetres, def:physics:phyx-mini-0444:phyxminiproblems-problemphyxmini0444-massinkilograms, def:physics:phyx-mini-0444:phyxminiproblems-problemphyxmini0444-specificvolumeincubicmetresperkilogram, def:physics:phyx-mini-0444:phyxminiproblems-problemphyxmini0444-temperatureinkelvin, def:physics:phyx-mini-0444:phyxminiproblems-problemphyxmini0444-twotankr410aexperiment, def:physics:phyx-mini-0444:phyxminiproblems-problemphyxmini0444-phasemassinkilograms, def:physics:phyx-mini-0444:phyxminiproblems-problemphyxmini0444-totalmassinkilograms, def:physics:phyx-mini-0444:phyxminiproblems-problemphyxmini0444-phasevolumeincubicmetres, def:physics:phyx-mini-0444:phyxminiproblems-problemphyxmini0444-totaloccupiedvolumeincubicmetres}
  The governing endpoint physics: Celsius and kelvin readouts refer to the same absolute temperature; each saturated phase obeys V\_phase = m\_phase v\_phase; non-evacuated contents fill their rigid tank; their pressure is the saturation pressure at their temperature; the integrated vapor stream is exactly the mass lost by A and gained by initially evacuated B. These laws are uniform physical relations and contain no numerical quality
\end{definition}
\begin{proof}
  This is the declared model definition of Satisfies Isothermal Saturated Transfer Laws.
\end{proof}
\begin{definition}[Answer Choice]
  \label{def:physics:phyx-mini-0444:phyxminiproblems-problemphyxmini0444-answerchoice}
  \lean{PhyXMiniProblems.ProblemPhyXMini0444.AnswerChoice}
  Labels of the four printed answer choices
\end{definition}
\begin{proof}
  This is the declared model definition of Answer Choice.
\end{proof}
\begin{definition}[displayed Quality Change]
  \label{def:physics:phyx-mini-0444:phyxminiproblems-problemphyxmini0444-displayedqualitychange}
  \lean{PhyXMiniProblems.ProblemPhyXMini0444.displayedQualityChange}
  \uses{def:physics:phyx-mini-0444:phyxminiproblems-problemphyxmini0444-answerchoice}
  Dimensionless quality changes printed beside the answer labels
\end{definition}
\begin{proof}
  This is the declared model definition of displayed Quality Change.
\end{proof}
\begin{definition}[recorded Dataset Answer]
  \label{def:physics:phyx-mini-0444:phyxminiproblems-problemphyxmini0444-recordeddatasetanswer}
  \lean{PhyXMiniProblems.ProblemPhyXMini0444.recordedDatasetAnswer}
  \uses{def:physics:phyx-mini-0444:phyxminiproblems-problemphyxmini0444-answerchoice}
  The dataset's recorded label, retained only as answer metadata
\end{definition}
\begin{proof}
  This is the declared model definition of recorded Dataset Answer.
\end{proof}
\begin{definition}[Is Reported Quality Change]
  \label{def:physics:phyx-mini-0444:phyxminiproblems-problemphyxmini0444-isreportedqualitychange}
  \lean{PhyXMiniProblems.ProblemPhyXMini0444.IsReportedQualityChange}
  \uses{def:physics:phyx-mini-0444:phyxminiproblems-problemphyxmini0444-twotankr410aexperiment, def:physics:phyx-mini-0444:phyxminiproblems-problemphyxmini0444-qualitychangeintanka, def:physics:phyx-mini-0444:phyxminiproblems-problemphyxmini0444-answerchoice, def:physics:phyx-mini-0444:phyxminiproblems-problemphyxmini0444-displayedqualitychange}
  Agreement with a three-decimal display: the exact modeled change lies within half of one unit in the last printed decimal place
\end{definition}
\begin{proof}
  This is the declared model definition of Is Reported Quality Change.
\end{proof}
% --- Archon declaration topology end ---
% --- Archon physics formalization source end ---
